\setcounter{section}{17}






  \zwischenueberschrift{Übungsaufgaben}

\inputaufgabe
{}
{





Berechne

\mathdisp {\left(  4 T^{(-1,3)} -6 T^{(-2,5)} + 5T^{(0,-2)} +3 T^{(-1,1)}  \right) \cdot \left(  -7 T^{(1,0)} + 8 T^{(4,5)} -4 T^{(-3,5)} +6 T^{(3,-1)}  \right)} {  }

im
\definitionsverweis {Monoidring}{}{}

\mathl{K[\Z^2]}{.}





}
{} {}

\inputaufgabe
{}
{





Berechne

\mathdisp {\left(  4 T^{0} -9 T^{1} + 9T^{2} +7 T^{3} -4T^4  \right) \cdot \left(  4 T^{0} -  T^{1} + 5T^{2} +7 T^{3} -T^4  \right)} {  }

im
\definitionsverweis {Monoidring}{}{}

\mathl{K[ \Z/( 5 ) ]}{.}





}
{} {}

\inputaufgabegibtloesung
{}
{





Berechne

\mathdisp {\left(  3 T^{0} -2 T^{1} + 5 T^{2}   \right) \cdot \left(  4 T^{0} - 6 T^{1} + 5T^{2}    \right)} {  }

im
\definitionsverweis {Monoidring}{}{}

\mathl{\Z/(7) [  \Z/(3 ) ]}{} über dem Körper
\mathl{\Z/(7)}{.}





}
{} {}

\inputaufgabe
{}
{





Zeige, dass die Multiplikation auf einem
\definitionsverweis {Monoidring}{}{}
zu einem
\definitionsverweis {kommutativen Monoid}{}{}
das Assoziativgesetz und das Distributivgesetz erfüllt.





}
{} {}

\inputaufgabe
{}
{





Es sei $K$ ein
\definitionsverweis {Körper}{}{.}
Finde ein
\definitionsverweis {kommutatives Monoid}{}{}
$M$ derart, dass eine
\definitionsverweis {Isomorphie}{}{}

\mavergleichskettedisp
{\vergleichskette
{ K[M]
}
{ \cong} { K[X,Y,U,V]/(UX-VY)
}
{ } {
}
{ } {
}
{ } {
}
}
{}{}{}
vorliegt.





}
{} {}

\inputaufgabe
{}
{





Es sei $R$ ein kommutativer Ring. Beweise die
$R$-\definitionsverweis {Algebraiso\-mor\-phie}{}{}

\mavergleichskettedisp
{\vergleichskette
{ R[\Z^n]
}
{ \cong} { R[X_1 , \ldots ,  X_n]_{X_1 \cdots X_n }
}
{ } {
}
{ } {
}
{ } {
}
}
{}{}{}
mithilfe der universellen Eigenschaften von Monoidringen und Nenneraufnahmen.





}
{} {}

\inputaufgabe
{}
{





Zeige, dass man den Koordinatenring zum Standardkegel
\mathl{V \left(  Z^2-X^2-Y^2  \right)}{} über ${\mathbb C}$ als einen
\definitionsverweis {Monoidring}{}{}
realisieren kann.





}
{} {}

\inputaufgabe
{}
{





Es sei $R$ ein
\definitionsverweis {kommutativer Ring}{}{}
und

\mavergleichskette
{\vergleichskette
{ A
}
{ = }{ R[X_1 , \ldots , X_n]
}
{  }{
}
{    }{
}
{   }{
}
}
{}{}{}
der
\definitionsverweis {Polynomring}{}{}
über $R$. Es sei

\mathdisp {\left\{ X^\nu  \mid   \nu \in J \subseteq \N^n         \right\}} {  }

eine Familie von
\definitionsverweis {Monomen}{}{}
in $A$. Zeige, dass die durch diese Monome erzeugte
$R$-\definitionsverweis {Unteralgebra}{}{}
von $A$ mit dem
\definitionsverweis {Monoidring}{}{}
$R[M]$ übereinstimmt, wobei $M$ das von den

\mathbed {\nu} {}
 {\nu \in J} {}
 {} {} {} {,}
\definitionsverweis {erzeugte Untermonoid}{}{}
von $\N^n$ ist.





}
{} {}

\inputaufgabe
{}
{





Es sei $M$ ein kommutatives
\definitionsverweis {Monoid}{}{}
und es sei

\mavergleichskette
{\vergleichskette
{ e
}
{ \in }{ M
}
{  }{
}
{    }{
}
{   }{
}
}
{}{}{}
ein Element mit

\mavergleichskette
{\vergleichskette
{ 2e
}
{ = }{ e
}
{  }{
}
{    }{
}
{   }{
}
}
{}{}{.}
Es sei $K$ ein
\definitionsverweis {Körper}{}{}
und es sei
\mathl{K[M]}{} der zugehörige
\definitionsverweis {Monoidring}{}{.}
Zeige, dass $T^e$ ein
\definitionsverweis {idempotentes Element}{}{}
in
\mathl{K[M]}{} ist.





}
{} {}

\inputaufgabe
{}
{





Es sei $G$ eine endliche
\definitionsverweis {Gruppe}{}{}

\mavergleichskette
{\vergleichskette
{G
}
{ \neq }{ 0
}
{  }{
}
{    }{
}
{   }{
}
}
{}{}{.}
Zeige, dass der
\definitionsverweis {Monoidring}{}{}

\mathl{{\mathbb C}[G]}{} nicht
\definitionsverweis {zusammenhängend}{}{}
ist, obwohl es in der Gruppe außer $0$ kein Element $e$ gibt, das die Gleichung

\mavergleichskette
{\vergleichskette
{ 2e
}
{ = }{ e
}
{  }{
}
{    }{
}
{   }{
}
}
{}{}{}
erfüllt.





}
{} {}

\inputaufgabe
{}
{





Man gebe ein Beispiel eines
\definitionsverweis {Untermonoids}{}{}

\mavergleichskette
{\vergleichskette
{ M
}
{ \subseteq }{ \N^2
}
{  }{
}
{    }{
}
{   }{
}
}
{}{}{,}
das nicht
\definitionsverweis {endlich erzeugt}{}{}
ist.





}
{} {}

Die invertierbaren Elemente in einem Monoid nennt man auch \definitionswort {Einheiten}{} des Monoids. Sie bilden die \definitionswort {Einheitengruppe}{} des Monoids.

\inputaufgabegibtloesung
{}
{





Es sei $M$ ein
\definitionsverweis {kommutatives Monoid}{}{}
und $K$ ein
\definitionsverweis {Körper}{}{.}
Es sei

\mavergleichskette
{\vergleichskette
{ m
}
{ \in }{ M
}
{  }{
}
{    }{
}
{   }{
}
}
{}{}{}
und

\mavergleichskette
{\vergleichskette
{ T^m
}
{ \in }{ K[M]
}
{  }{
}
{    }{
}
{   }{
}
}
{}{}{.}
Zeige, dass $m$ genau dann eine
\definitionsverweis {Einheit}{}{}
in $M$ ist, wenn $T^m$ eine
\definitionsverweis {Einheit}{}{}
in
\mathl{K[M]}{} ist.





}
{} {}

\inputaufgabe
{}
{





Zeige, dass die
\definitionsverweis {Differenzengruppe}{}{}
zu einem
\definitionsverweis {kommutativen Monoid}{}{}
in der Tat eine
\definitionsverweis {Gruppe}{}{}
ist.





}
{} {}

\inputaufgabe
{}
{





Es sei $M$ ein kommutatives
\definitionsverweis {Monoid}{}{.}
Zeige, dass die zugehörige
\definitionsverweis {Differenzengruppe}{}{}

\mavergleichskette
{\vergleichskette
{ \Gamma
}
{ = }{ \Gamma(M)
}
{  }{
}
{    }{
}
{   }{
}
}
{}{}{}
eine kommutative
\definitionsverweis {Gruppe}{}{}
ist, und dass sie folgende universelle Eigenschaft besitzt: Zu jedem
\definitionsverweis {Monoidhomomorphismus}{}{}
\maabbdisp {\varphi} { M } { G
} {}
in eine Gruppe $G$ gibt es einen eindeutig bestimmten
\definitionsverweis {Gruppenhomomorphismus}{}{}
\maabbdisp {\tilde{\varphi}} {\Gamma} {G
} {,}
der $\varphi$ fortsetzt.





}
{} {}

\inputaufgabe
{}
{





Es seien

\mavergleichskette
{\vergleichskette
{ M
}
{ \subseteq }{ N
}
{  }{
}
{    }{
}
{   }{
}
}
{}{}{}
\definitionsverweis {kommutative Monoide}{}{.}
Zeige, dass durch

\mavergleichskettedisp
{\vergleichskette
{ \tilde{M}
}
{ =} { \left\{ n \in N  \mid  \text{es gibt } k \in \N_+ \text{ mit } kn \in M        \right\}
}
{ } {
}
{ } {
}
{ } {
}
}
{}{}{}
ein Untermonoid von $N$ gegeben ist, das $M$ umfasst.





}
{} {}

\inputaufgabe
{}
{





Es sei $M$ ein kommutatives
\definitionsverweis {Monoid}{}{}
mit zugehöriger
\definitionsverweis {Differenzengruppe}{}{}

\mavergleichskette
{\vergleichskette
{ \Gamma
}
{ = }{ \Gamma(M)
}
{  }{
}
{    }{
}
{   }{
}
}
{}{}{.}
Zeige, dass folgende Aussagen äquivalent sind.
\aufzaehlungdrei{ $M$ ist ein
\definitionsverweis {Monoid mit Kürzungsregel}{}{.}
}{Die kanonische Abbildung
\maabb {} { M } {\Gamma(M)} {}
ist injektiv.
}{ $M$ lässt sich als Untermonoid einer Gruppe realisieren.
}





}
{} {}

\inputaufgabe
{}
{





Es sei $M$ ein
\definitionsverweis {kommutatives Monoid}{}{}
und $R$ ein
\definitionsverweis {kommutativer Ring}{}{.}
Charakterisiere, für welche Teilmengen

\mavergleichskette
{\vergleichskette
{ I
}
{ \subseteq }{ M
}
{  }{
}
{    }{
}
{   }{
}
}
{}{}{}
die Teilmenge

\mavergleichskettedisp
{\vergleichskette
{ R[I]
}
{ =} { \bigoplus_{m \in I} T^m
}
{ \subseteq} { R[M]
}
{ } {
}
{ } {
}
}
{}{}{}
ein
\definitionsverweis {Ideal}{}{}
in $R[M]$ ist.





}
{} {}

\inputaufgabe
{}
{





Betrachte den Monoidhomomorphismus

\mathdisp {\N^2 \longrightarrow \Z,\, e_1 \longmapsto 1,\, e_2 \longmapsto -1} { . }

Beschreibe die zugehörige Abbildung zwischen den Monoidringen (für einen Körper $K$) und den zugehörigen $K$-Spektren.





}
{} {}

\inputaufgabe
{}
{





Wir betrachten die
\definitionsverweis {kommutativen}{}{}
\definitionsverweis {Monoide}{}{}

\mavergleichskette
{\vergleichskette
{ M
}
{ = }{ \N^r
}
{  }{
}
{    }{
}
{   }{
}
}
{}{}{}
und

\mavergleichskette
{\vergleichskette
{ N
}
{ = }{ \N^s
}
{  }{
}
{    }{
}
{   }{
}
}
{}{}{.}
Zeige, dass ein
\definitionsverweis {Monoidhomomorphismus}{}{}
von $M$ nach $N$ eindeutig durch eine
\definitionsverweis {Matrix}{}{}
\zusatzklammer {mit $r$ Spalten und $s$ Zeilen} {} {}
mit Einträgen aus $\N$ bestimmt ist.





}
{Wie sieht die zugehörige Spektrumsabbildung aus?} {}

\inputaufgabe
{}
{





Es sei $M$ eine quadratische $r\times r$-Matrix mit Einträgen aus $\N$ mit der zugehörigen Monoidabbildung und der zugehörigen Spektrumsabbildung
\maabbdisp {\varphi} { K\!-\!\operatorname{Spek}\, \left(  K[\N^r]  \right) } { K\!-\!\operatorname{Spek}\, \left(  K[\N^r]  \right)
} {,}
wobei $K$ ein unendlicher Körper sei. Zeige, dass  genau dann

\mavergleichskettedisp
{\vergleichskette
{ \det  M
}
{ \neq} { 0
}
{ } {
}
{ } {
}
{ } {
}
}
{}{}{}
ist, wenn $\varphi$ surjektiv auf eine offene Menge aus $K\!-\!\operatorname{Spek}\, \left(  K[\N^r]  \right)$ abbildet.





}
{} {}

Die nächste Aufgabe verwendet die folgende Definition.





Ein \definitionswort {Filter}{} $F$ in einem
\definitionsverweis {kommutativen Monoid}{}{}
$M$ ist ein
\definitionsverweis {Untermonoid}{}{,}
das zusätzlich \definitionswort {teilerstabil}{} ist. D.h. falls $f \in F$ ist und $g|f$ gilt, so ist auch $g \in F$.







\inputaufgabe
{}
{





Es sei $M$ ein
\definitionsverweis {kommutatives Monoid}{}{.}
Zeige, dass es in $M$ einen kleinsten
\definitionsverweis {Filter}{}{}
gibt und dass dieser eine
\definitionsverweis {Gruppe}{}{}
bildet.





}
{} {}

\inputaufgabe
{}
{





Es sei $M$ ein kommutatives
\definitionsverweis {Monoid}{}{}
und sei

\mavergleichskette
{\vergleichskette
{ f
}
{ \in }{ M
}
{  }{
}
{    }{
}
{   }{
}
}
{}{}{.}
Wir betrachten die Menge

\mavergleichskettedisp
{\vergleichskette
{ M_f
}
{ =} { \left\{  m-nf  \mid   m \in M , \,  n \in \N       \right\}  /\sim
}
{ } {
}
{ } {
}
{ } {
}
}
{}{}{,}
wobei die Relation

\mavergleichskettedisp
{\vergleichskette
{ m-nf
}
{ \sim} { m' -n'f
}
{ } {
}
{ } {
}
{ } {
}
}
{}{}{}
genau dann gilt, wenn es ein

\mavergleichskette
{\vergleichskette
{ k
}
{ \in }{ \N
}
{  }{
}
{    }{
}
{   }{
}
}
{}{}{}
derart gibt, dass

\mavergleichskettedisp
{\vergleichskette
{ m +n'f + kf
}
{ =} { m' +nf + kf
}
{ } {
}
{ } {
}
{ } {
}
}
{}{}{}
in $M$ gilt.
\aufzaehlungdreiabc{Zeige, dass $\sim$ eine
\definitionsverweis {Äquivalenzrelation}{}{}
ist.
}{Definiere auf $M_f$ eine Monoidstruktur.
}{Es sei $R$ ein
\definitionsverweis {kommutativer Ring}{}{}
und sei

\mavergleichskette
{\vergleichskette
{ T^f
}
{ \in }{ R[M]
}
{  }{
}
{    }{
}
{   }{
}
}
{}{}{}
das Monom zu $f$ im Monoidring. Zeige

\mavergleichskettedisp
{\vergleichskette
{ R[M_f ]
}
{ \cong} { R[M]_{T^f}
}
{ } {
}
{ } {
}
{ } {
}
}
{}{}{.}
}





}
{} {}

\inputaufgabe
{}
{





Es seien $M,N$ endlich erzeugte kommutative Monoide mit den $K$-Spektren
\mathl{K\!-\!\operatorname{Spek}\, \left(  K[M]  \right)= \operatorname{Mor}_{ \operatorname{mon} } \, ( M ,  K)}{} und
\mathl{K\!-\!\operatorname{Spek}\, \left(  K[N]  \right)= \operatorname{Mor}_{ \operatorname{mon} } \, ( N ,  K)}{.} Zeige, dass man für einen Monoidhomomorphismus
\mathl{\varphi:M \rightarrow N}{} die zugehörige Spektrumsabbildung auf zwei verschiedene Weisen definieren kann, die aber inhaltlich übereinstimmen.





}
{} {}

\inputaufgabe
{}
{





Wir betrachten das kommutative
\definitionsverweis {Monoid}{}{}
$M$, das durch die drei Erzeuger
\mathl{e,f,g}{} und die einzige Relation

\mavergleichskette
{\vergleichskette
{e+f
}
{ = }{ 5g
}
{  }{
}
{    }{
}
{   }{
}
}
{}{}{}
gegeben ist. Bestimme das
$K$-\definitionsverweis {Spektrum}{}{}
zu $M$ für verschiedene
\definitionsverweis {Körper}{}{}
$K$.





}
{} {}

\inputaufgabe
{}
{





Es sei $M$ ein endliches kommutatives
\definitionsverweis {Monoid}{}{.}
Zeige, dass das
$K$-\definitionsverweis {Spektrum}{}{}
zu $M$ auch endlich ist.





}
{} {}

\inputaufgabe
{}
{





Berechne in

\mavergleichskette
{\vergleichskette
{ R
}
{ = }{ \R[\Q_{\geq 0}]
}
{  }{
}
{    }{
}
{   }{
}
}
{}{}{}
das Produkt

\mathdisp {\left(  X^2 +4 X^{3/2}-5X+ X^{1/2}   \right) \left(  2 X^{3/2}+4X-7X^{1/2}   \right)} { . }






}
{} {}

\inputaufgabe
{}
{





Zeige, dass man jedes Element

\mavergleichskette
{\vergleichskette
{ F
}
{ \in }{ R
}
{ = }{ K[\Q_{\geq 0}]
}
{    }{
}
{   }{
}
}
{}{}{}
\zusatzklammer {$K$ ein Körper} {} {}
als ein Polynom in
\mathl{X^{1/b}}{} mit einem

\mavergleichskette
{\vergleichskette
{ b
}
{ \in }{ \N_+
}
{  }{
}
{    }{
}
{   }{
}
}
{}{}{}
schreiben kann, dass es also ein

\mavergleichskette
{\vergleichskette
{ P
}
{ \in }{ K[Y]
}
{  }{
}
{    }{
}
{   }{
}
}
{}{}{}
derart gibt, dass

\mavergleichskette
{\vergleichskette
{ F
}
{ = }{ P(X^{1/b})
}
{  }{
}
{    }{
}
{   }{
}
}
{}{}{}
gilt. Welches Polynom kann man bei

\mavergleichskettedisp
{\vergleichskette
{ F
}
{ =} { X^{1/2} +X^{1/3} + X^{1/5}
}
{ } {
}
{ } {
}
{ } {
}
}
{}{}{}
nehmen?





}
{} {}

\inputaufgabe
{}
{





Zeige, dass in

\mavergleichskette
{\vergleichskette
{ R
}
{ = }{ K[\Q_{\geq 0}]
}
{  }{
}
{    }{
}
{   }{
}
}
{}{}{}
das Element $X$ keine Zerlegung in irreduzible Elemente besitzt.





}
{} {}

\inputaufgabe
{}
{





Zeige, dass in

\mavergleichskette
{\vergleichskette
{ R
}
{ = }{ \R [\Q_{\geq 0}]
}
{  }{
}
{    }{
}
{   }{
}
}
{}{}{}
das Element $X^2+1$ nicht
\definitionsverweis {irreduzibel}{}{}
ist.





}
{} {}

\inputaufgabe
{}
{





Zeige, dass es in

\mavergleichskette
{\vergleichskette
{ R
}
{ = }{ {\mathbb C}[\Q_{\geq 0}]
}
{  }{
}
{    }{
}
{   }{
}
}
{}{}{}
keine
\definitionsverweis {irreduziblen Elemente}{}{}
gibt.





}
{} {}

\inputaufgabegibtloesung
{}
{





Bestimme sämtliche Teiler von $X$ im
\definitionsverweis {Ring}{}{}

\mavergleichskette
{\vergleichskette
{ R
}
{ = }{ K[\Q_{\geq 0}]
}
{  }{
}
{    }{
}
{   }{
}
}
{}{}{,}
wobei $K$ ein
\definitionsverweis {Körper}{}{}
ist.





}
{} {}

\inputaufgabegibtloesung
{}
{





Bestimme die
\definitionsverweis {Einheiten}{}{}
im Ring

\mavergleichskette
{\vergleichskette
{ R
}
{ = }{ K[\Q]
}
{  }{
}
{    }{
}
{   }{
}
}
{}{}{,}
wobei $K$ ein
\definitionsverweis {Körper}{}{}
ist.





}
{} {}





  \zwischenueberschrift{Aufgaben zum Abgeben}

\inputaufgabe
{6}
{





Es seien

\mavergleichskette
{\vergleichskette
{ M
}
{ \subseteq }{ N
}
{  }{
}
{    }{
}
{   }{
}
}
{}{}{}
endlich erzeugte kommutative Monoide mit Kürzungsregel. Zeige, dass für einen
\definitionsverweis {Körper}{}{}
$K$ der
\definitionsverweis {Ringhomomorphismus}{}{}

\mavergleichskette
{\vergleichskette
{ K[M]
}
{ \subseteq }{ K[N]
}
{  }{
}
{    }{
}
{   }{
}
}
{}{}{}
genau dann
\definitionsverweis {endlich}{}{}
ist, wenn es zu jedem

\mavergleichskette
{\vergleichskette
{ n
}
{ \in }{ N
}
{  }{
}
{    }{
}
{   }{
}
}
{}{}{}
ein

\mavergleichskette
{\vergleichskette
{ k
}
{ \in }{ \N_+
}
{  }{
}
{    }{
}
{   }{
}
}
{}{}{}
mit

\mavergleichskette
{\vergleichskette
{ kn
}
{ \in }{ M
}
{  }{
}
{    }{
}
{   }{
}
}
{}{}{}
gibt.





}
{} {}

\inputaufgabe
{4}
{





Es sei

\mavergleichskette
{\vergleichskette
{ M
}
{ = }{ (\Q,+)
}
{  }{
}
{    }{
}
{   }{
}
}
{}{}{}
die additive Gruppe der
\definitionsverweis {rationalen Zahlen}{}{.}
Bestimme
\mathl{\Q\!-\!\operatorname{Spek}\, \left(  \Q[M]   \right)}{.} Wie sieht es aus, wenn man $\Q$ durch $\R$ ersetzt?





}
{} {}

\inputaufgabe
{4}
{





Es sei
\maabb {\varphi} { M } { N
} {}
ein
\definitionsverweis {Homomorphismus}{}{}
von kommutativen Monoiden. Zeige, dass die Menge aller Punkte aus
\mathl{K\!-\!\operatorname{Spek}\, \left(   K[N]   \right)}{,} die unter der Spektrumsabbildung auf den Einspunkt

\mavergleichskette
{\vergleichskette
{ 1
}
{ \in }{ K\!-\!\operatorname{Spek}\, \left(   K[M]   \right)
}
{  }{
}
{    }{
}
{   }{
}
}
{}{}{}
\zusatzklammer {das ist der Punkt, der der konstanten Abbildung $M \mapsto 1$ entspricht} {} {} abgebildet werden, selbst die Struktur eines $K$-Spektrums eines geeigneten Monoids besitzt.





}
{} {}

\inputaufgabe
{4}
{





Wir betrachten Monoide der Form

\mavergleichskettedisp
{\vergleichskette
{ M
}
{ =} { ( \Z/( m ) ,+)
}
{ } {
}
{ } {
}
{ } {
}
}
{}{}{}
Beschreibe
\mathl{K\!-\!\operatorname{Spek}\, \left(   K[M]   \right)}{} allgemein sowie für die
\definitionsverweis {Körper}{}{}

\mavergleichskette
{\vergleichskette
{ K
}
{ = }{ \R, {\mathbb C}, \Z/( 5 )
}
{  }{
}
{    }{
}
{   }{
}
}
{}{}{.}
Finde die
\definitionsverweis {idempotenten Elemente}{}{}
von
\mathl{{\mathbb C} [ \Z/( 3 ) ]}{.}





}
{} {}

\inputaufgabe
{4}
{





Es sei $M$ ein kommutatives Monoid. Definiere eine Bijektion zwischen den folgenden Objekten.
\aufzaehlungvier{ \definitionsverweis {Filter}{}{}
in $M$.
}{
\mathl{\operatorname{Mor}_{ \operatorname{mon} } \, ( M , (\{0,1\},1,\cdot))}{.}
}{
\mathl{{\mathbb F}_2-\operatorname{Spek}  \,  (M)}{}
}{
\mathl{\left\{ \varphi \in K\!-\!\operatorname{Spek}\, \left(  K[M]  \right)   \mid   \varphi(M) \subseteq \{0,1\}         \right\}}{.}
\zusatzklammer {Dabei ist $K$ ein Körper.} {} {}
}





}
{} {}

\inputaufgabe
{3}
{





Es seien $M$ und $N$
\definitionsverweis {kommutative Monoide}{}{}
und sei $K$ ein
\definitionsverweis {Körper}{}{.}
In welcher Beziehung steht
\mathl{K\!-\!\operatorname{Spek}\, \left(  K[M \times N]  \right)}{} zu
\mathl{K\!-\!\operatorname{Spek}\, \left(  K[M]  \right)}{} und
\mathl{K\!-\!\operatorname{Spek}\, \left(  K[N]  \right)}{?}





}
{} {}

\inputaufgabe
{4}
{





Es sei $K$ ein
\definitionsverweis {Körper}{}{}
und $G$ eine
\definitionsverweis {Gruppe}{}{.}
Dann können wir den
\definitionsverweis {Monoidring}{}{}
$K[G]$ betrachten. Es sei nun weiter $M$ ein $K[G]$-Modul. Zeige, dass
\aufzaehlungzweiabc{ $M$ nichts anderes ist als ein $K$-Vektorraum $V$ zusammen mit einem
\definitionsverweis {Gruppenhomomorphismus}{}{}
\maabb {\rho} { G} { \operatorname{Aut}_K(V)
} {.}
}{ein $K[G]$-Modulhomomorphismus
\maabb {\varphi} { M } { M
} {}
eine $K$-lineare Abbildung ist, für die zusätzlich

\mavergleichskette
{\vergleichskette
{ \varphi \circ \rho(g)
}
{ = }{ \rho \circ \varphi
}
{  }{
}
{    }{
}
{   }{
}
}
{}{}{}
für alle

\mavergleichskette
{\vergleichskette
{ g
}
{ \in }{ G
}
{  }{
}
{    }{
}
{   }{
}
}
{}{}{}
gilt.
}





}
{Bemerkung: $\rho$ heißt dann eine \stichwort {Darstellung} {} von $G$. Solche Darstellungen sind oft einfacher zu handhaben als $G$ und man kann mit Hilfe von $\rho$ oft hilfreiche Erkenntnisse über $G$ selbst gewinnen.} {}
